This section introduces the DCC model and discusses its maximum likelihood estimation in large dimensions. In the following discussion and throughout the thesis, the subscript $i$ indexes the variables and runs from $1$ to $N$, and the subscript $t$ indexes time and covers the range of integers from $1$ to $T$.

\subsection{Model Definition}
Let $r_{it}$ denote the log return on asset $i$ at time $t$ and let $X_t = (r_{1t}, r_{2t}, \dots , r_{Nt})^\prime$ be the vector of log returns containing all observations on the $N$ assets at time $t$. Let also $\mathcal{F}_{t}$ denote the information set generated by the history of returns up to time $t$. We make the simplifying assumption that $\mathbb{E}[X_t \vert{} \mathcal{F}_{t-1}] = 0$, which is common and reasonable for daily returns. Let $Z_t$ be an $N \times 1$ vector of independently and identically distributed (i.i.d.) random variables with zero mean and unit variance. Then, in the most general and abstract terms, the stochastic process for $X_t$ is defined as:
\begin{align}
	X_t = H_t^{1/2} Z_t,
\end{align}
where $H_t = \operatorname{Cov}(X_t \vert{} \mathcal{F}_{t-1})$ is the conditional covariance matrix of $X_t$. The problem then becomes specifying a model for $H_t$ and estimating its parameters from sample data. The literature on modeling the matrix $H_t$ through parametric models has a rich and long history (see \textcite{BAUWENS} for a survey), and the typical approach consists of specifying a time-series model for $H_t$ and estimating its parameters via maximum likelihood. Among the many possible alternatives that exist, we focus on the DCC model of \textcite{ENGLE2002} for its simplicity and scalability. The DCC model is based on the extension of univariate GARCH models to the multivariate case. The main approach is as follows. First, a univariate GARCH model is defined for each of the return series. We focus here on the baseline GARCH(1,1) model specification. Therefore, for each $r_{it}$ we have the following specification:
\begin{align}
	r_{it} = d_{it}(\mathcal{F}_{t-1};\theta) z_{it},
\end{align}
with $z_{it}$ a normally and serially independent random variable with zero mean and unit variance that is orthogonal to $r_{it}$. The GARCH(1,1) model for $d_{it}$ is defined as:
\begin{align}
	d_{it}^2 = \omega_i + a_i r_{it-1}^2 + b_i d_{it-1}^2.
\end{align}
With the conditional variance of the return series estimated via the process defined above, we can obtain the devolatilized return series as $s_{it} = r_{it}/d_{it}$ for each asset, which enter directly into the dynamic correlation equation in the DCC model. Let $S_{t} = (s_{1t}, s_{2t}, \dots , s_{Nt})^\prime$ and let $D_t$ be the $N$-dimensional diagonal matrix whose $i$th diagonal element is $d_{it}$. Furthermore, let $C = \operatorname{Corr}(X_t) = \operatorname{Cov}(S_t)$ be the unconditional correlation matrix. Then the DCC model is defined as follows:
\begin{align}
	Q_t &= (1 - \alpha - \beta)C + \alpha S_{t-1} S_{t-1}^\prime + \beta Q_{t-1}, \label{eq:sec2_dcc_q} \\
	R_t &= \operatorname{diag}(Q_t)^{-1/2} Q_t \operatorname{diag}(Q_t)^{-1/2}, \label{eq:sec2_dcc_r} \\
	H_t &= D_t R_t D_t. \label{eq:sec2_dcc_h}
\end{align}

\subsection{Maximum Likelihood Estimation}

Assuming that the underlying data generating process is driven by the multivariate normal law, the log-likelihood function for the DCC model can be expressed up to a constant as:
\begin{align}
	&X_t | \mathcal{F}_{t-1} \sim \mathcal{N}(0,H_t)\, \\
	&L = -\frac{1}{2} \sum_{t=1}^T (\operatorname{log} \left( \left| H_t \right| \right) + X_t^\prime H_t^{-1} X_t), \label{eq:sec2_log_like_1}
\end{align}
Where $|\cdot|$ denotes the determinant. Substituting $H_t$ with the expression \eqref{eq:sec2_dcc_h} in the log likelihood function yields:
\begin{align}
	L = -\frac{1}{2} \sum_{t=1}^T (\operatorname{log}(|D_tR_tD_t|) + X_t^\prime D_t^{-1}R_t^{-1}D_t^{-1}X_t )
\end{align}

Since by definition we have $S_t = D_t^{-1/2}X_t$, we can rewrite the log likelihood function as:
\begin{align}
	L = -\frac{1}{2} \sum_{t=1}^T (2\operatorname{log}(|D_t|) + X_t^\prime D_t^{-1}D_t^{-1}X_t - S_t^\prime S_t +  \operatorname{log}(|R_t|) + S_t^\prime R_t^{-1}S_t ), \label{eq:sec2_log_like_2}
\end{align}
where we used the multiplicative property of the determinant and added and subtracted the two equivalent terms $X_t^\prime D_t^{-1}D_t^{-1}X_t$ and $S_t^\prime S_t$. The expression in \eqref{eq:sec2_log_like_2} allows for the natural decomposition of the log likelihood function into the sum of a volatility component and a correlation component. Let the univariate GARCH parameters in $D_t$ be denoted $\theta$ and the additional parameters in $R_t$ be denoted $\phi$. Then the log likelihood function in \eqref{eq:sec2_log_like_2} can be written as:
\begin{align}
	L(\theta,\phi) = L_V(\theta) + L_C(\theta,\phi), \label{eq:sec2_log_like_decomp}
\end{align}
where
\begin{align}
	L_V(\theta) = -\frac{1}{2} \sum_{t=1}^T (2 \operatorname{log}(|D_t|) + X_t^\prime D_t^{-1}D_t^{-1}X_t), \label{eq:sec2_log_like_v}
\end{align}
is the volatility term and
\begin{align}
	L_C(\theta, \phi) = -\frac{1}{2}\sum_{t=1}^T (\operatorname{log}(|R_t|) + S_t^\prime R_t^{-1}S_t - S_t^\prime S_t), \label{eq:sec2_log_like_c}
\end{align}
is the correlation term. The main strength of the DCC model is that it allows for the univariate GARCH parameters to be estimated separately from the conditional correlation parameters. The volatility term $L_V(\theta)$ is just the sum of the individual GARCH log likelihoods:
\begin{align}
	L_V(\theta) &= -\frac{1}{2} \sum_{t=1}^T (\operatorname{log}(|D_t|) + X_t^\prime D_t^{-1} D-t^{-1}X_t), \\
	 &= -\frac{1}{2} \sum_{t=1}^T \left(\operatorname{log}\left( \prod_{i=1}^N d_{it} \right) + S_t^\prime S_t \right), \\
	 &= -\frac{1}{2} \sum_{t=1}^T \left( \sum_{i=1}^N \operatorname{log}(d_{it}) + s_{it}^2 \right), \\ 
	 &= - \frac{1}{2} \sum_{t=1}^T\sum_{i=1}^N \left(\operatorname{log}(d_{it}) + \frac{r_{it}^2}{d_{it}^2}\right),
\end{align}

which is maximized by separately maximizing each term. Once estimates of the univariate GARCH parameters are obtained, we can run the recursions in equations \eqref{eq:sec2_dcc_q} and \eqref{eq:sec2_dcc_r}, and estimate the dynamic correlation parameters by maximizing the correlation component of the log likelihood function in \eqref{eq:sec2_log_like_c}. Therefore, the two-step approach to estimating the parameters of the DCC model via maximizing the joint log likelihood function in \eqref{eq:sec2_log_like_decomp} is to find:
\begin{align}
	\hat{\theta} = \operatorname{arg max}\{ L_V(\theta) \},
\end{align}
and next, taking this value as given in the second stage, maximize the correlation component over the correlation parameter space:
\begin{align}
	\hat{\phi} = \operatorname{arg max}_{\phi} \{ L_C(\hat{\theta}, \phi) \}.
\end{align}

\subsection{Modeling Issues in Large Dimensions}

Maximizing the correlation component $L_C(\theta, \phi)$ additionally requires an estimate of the unconditional correlation matrix $C$ in equation \eqref{eq:sec2_dcc_q}. When working in large dimensions, the usual estimator, the sample correlation matrix defined as $\hat{C} = T^{-1}\sum_{t=1}^T S_t S_t^\prime$, is no longer feasible since its eigenvalues are substantially biased and the matrix itself is ill-conditioned. We will turn to this issue in Section \ref{sec:shrink} where we explore the shrinkage estimation of the unconditional correlation matrix.

A second issue concerning the maximization of $L_C(\theta, \phi)$ in large dimensions is that the computation cost of doing so can be substantially high. Solving for the maximum likelihood estimates requires inverting the conditional correlation matrix $R_t$ in each iteration, which costs $O(N^3)$ operations. Therefore, the total computation cost of only inverting this matrix would be $O(N^3 \times T)$ over the full sample. Furthermore, inverting a large dimensional matrix such as $R_t$ can dramatically increase estimation error if the matrix is ill-conditioned. 

\textcite{PAKELETAL} propose a method for the estimation of multivariate GARCH models which doe not suffer from the computational and statistical issues outlined above. Their method is based on the composite likelihood function which approximates the full-dimensional joint likelihood function with combinations of lower dimensional bivariate densities. The advantage of their method is that it requires only the inversion of a $2 \times 2$ matrix in each iteration for which an analytical formula exists. Therefore the computational cost of this method is substantially lower than the computational cost of full information maximum likelihood estimation. The only drawback is that their method loses efficiency since it ignores the cross-correlations between the variables outside each pair when computing the marginal bivariate densities.
 
\subsection{Composite Likelihood Method}

The composite likelihood function is defined as the product of pairs of bivariate likelihood functions of the variables. For a total of $N$ variables, there are $\binom{N}{2}$ ways that the variables can be paired. For each pair of variables $r_i,r_j$ with $i \neq j$, the bivariate likelihood function is given up to a constant by:
\begin{align}
	f(r_{i1},r_{i2},\dots,r_{iT},r_{j1},r_{j2},\dots,r_{jT};\theta_i,\theta_j,\phi_{ij}) =  \prod_{t=1}^T |D_{ij,t}R_{ij,t}D_{ij,t}|^{-1/2} e^{-\frac{(r_{it}\, r_{jt})^\prime D_{ij,t}^{-1}R_{ij,t}^{-1}D_{ij,t}^{-1}(r_{it}\, r_{jt})}{2}},
\end{align}

where $D_{ij,t}$ is the diagonal matrix whose diagonal entries are given by $d_{it}$ and $d_{jt}$, and $R_{ij,t}$ is the $2 \times 2$ matrix whose diagonal entries are given by $i$th and $j$th diagonal entries, and its off-diagonal entries given by the $ij$th entry of the conditional correlation matrix in \eqref{eq:sec2_dcc_r}. The log-likelihood function for the pair $i,j$ is expressed as:
\begin{align}
	L_{ij} = -\frac{1}{2} \sum_{t=1}^T \left( 2\operatorname{log}(|D_{ij,t}|) + (r_{it}\, r_{jt})^\prime D_{ij,t}^{-1}D_{ij,t}^{-1} (r_{it}\, r_{jt}) - S_{ij,t}^\prime S_{ij,t} + \operatorname{log}(|R_{ij,t}|) + S_{ij,t}^\prime R_t^{-1} S_{ij,t} \right), \label{eq:sec2_comp_log_like_1}
\end{align}

where $S_{ij,t} = (r_{it}/d_{it}\, r_{jt}/d_{jt})$ is the $2 \times 1$ vector of devolatilized returns. Just as in the case of the full-dimensional likelihood, the log-likelihood function in \eqref{eq:sec2_comp_log_like_1} can be decomposed into a volatility part and a correlation part:
\begin{align}
	L_{ij} = L_{V,ij}(\theta) + L_{C,ij}(\theta, \phi)
\end{align}

The volatility term is just the sum of the individual GARCH log-likelihoods and is given by:

\begin{align}
	L_{V,ij}(\theta) = -\frac{1}{2} \sum_{t=1}^T \sum_{k \in [i,j]} \left( \operatorname{log}(d_{kt}) + \frac{r_{kt}^2}{d_{kt}^2} \right).
\end{align}

The correlation component is given like before as:
\begin{align}
	L_{C,ij}(\theta, \phi) = \sum_{t=1}^T ( \operatorname{log}(|R_{ij,t}|) + S_{ij,t}^\prime R_{ij,t}^{-1}S_{ij,t} - S_{ij,t}^\prime S_{ij,t}). 
\end{align}

Therefore, the same two-step logic in estimation of the parameters of the DCC model still applies. In order to obtain the composite log-likelihood function, we must take the average of all bivariate log-likelihoods. Depending on how the pairs are combined, we might have different composite log-likelihood functions. One way to pair the variables is to consider all possible combinations, with replacement, of pairs of variables. This will result in $\binom{N}{2} = N!/(N-2)!\times 2$ total pairs. Another way is to only consider the contiguous pairs of variables, that is, for $i=1,2,\dots,N$, the first pair consists of the variables $1,2$, the second pair of the variables $3,4$ and so on so forth until the last pair that consists of variables $N-1,N$. The difference between the two methods is that in the first, all possible pairs of assets are considered and the resulting composite log likelihood function is pooled over a broader set of joint bivaraite likelihoods. This way fewer information will be lost. In the second method, on the other hand, only the contiguous pairs of variables are considered and the maximum number of pairs included is $N-1$. Another difference between the two methods is that the second method is more scalable since there are $\binom{N}{2}/(N-1)$ distinct ways to choose $N-1$ contiguous pairs whereas there is only one way to choose all the $\binom{N}{2}$ possible pairs. Once the pairs have been chosen and the bivariate log likelihoods computed, the composite log likelihood function is constructed as:
\begin{align}
	L_{T,L}(\theta, \phi) = \frac{1}{L} \sum_{j=1}^L L_{j}(\theta, \phi), \label{eq:sec2_comp_like_2}
\end{align}

where $L$ is the number of pairs, and $L_{j}(\theta, \phi)$ is the bivaraite log likelihood function of the $j$th pair. The maximum composite likelihood estimator is then obtained by maximizing the composite log likelihood function in \eqref{eq:sec2_comp_like_2} in the two step procedure discussed earlier as:
\begin{align}
	\hat{\theta} = \text{arg max}\{ L_{V,TL}(\theta) \}, \\
	\hat{\phi} = \text{arg max}_{\phi}\{ L_{C,TL}(\hat{\theta},\phi) \}
\end{align}



















